
\documentclass{conm-p-l}
%%%%%%%%%%%%%%%%%%%%%%%%%%%%%%%%%%%%%%%%%%%%%%%%%%%%%%%%%%%%%%%%%%%%%%%%%%%%%%%%%%%%%%%%%%%%%%%%%%%%%%%%%%%%%%%%%%%%%%%%%%%%%%%%%%%%%%%%%%%%%%%%%%%%%%%%%%%%%%%%%%%%%%%%%%%%%%%%%%%%%%%%%%%%%%%%%%%%%%%%%%%%%%%%%%%%%%%%%%%%%%%%%%%%%%%%%%%%%%%%%%%%%%%%%%%%
\usepackage{amsfonts}

\setcounter{MaxMatrixCols}{10}
%TCIDATA{OutputFilter=LATEX.DLL}
%TCIDATA{Version=5.50.0.2960}
%TCIDATA{<META NAME="SaveForMode" CONTENT="1">}
%TCIDATA{BibliographyScheme=Manual}
%TCIDATA{Created=Wednesday, November 18, 2020 08:23:23}
%TCIDATA{LastRevised=Thursday, November 19, 2020 13:01:37}
%TCIDATA{<META NAME="GraphicsSave" CONTENT="32">}
%TCIDATA{<META NAME="DocumentShell" CONTENT="Author Packages for AMS\Contemporary Mathematics (Proceedings)">}
%TCIDATA{CSTFile=conm-p-l.cst}

\newtheorem{theorem}{Theorem}[section]
\newtheorem{lemma}[theorem]{Lemma}
\theoremstyle{definition}
\newtheorem{definition}[theorem]{Definition}
\newtheorem{example}[theorem]{Example}
\newtheorem{xca}[theorem]{Exercise}
\theoremstyle{remark}
\newtheorem{remark}[theorem]{Remark}
\numberwithin{equation}{section}
\theoremstyle{plain}
\newtheorem{acknowledgement}{Acknowledgement}
\newtheorem{algorithm}{Algorithm}
\newtheorem{axiom}{Axiom}
\newtheorem{case}{Case}
\newtheorem{claim}{Claim}
\newtheorem{conclusion}{Conclusion}
\newtheorem{condition}{Condition}
\newtheorem{conjecture}{Conjecture}
\newtheorem{corollary}{Corollary}
\newtheorem{criterion}{Criterion}
\newtheorem{exercise}{Exercise}
\newtheorem{notation}{Notation}
\newtheorem{problem}{Problem}
\newtheorem{proposition}{Proposition}
\newtheorem{solution}{Solution}
\newtheorem{summary}{Summary}
\copyrightinfo{2001}{enter name of copyright holder}

\input{tcilatex}

\begin{document}
\title[Short Title]{Contemporary Mathematics (Proceedings)}
\author{Author One}
\address[A. One and A. Two]{Author OneTwo address line 1\\
Author OneTwo address line 2}
\email[A. One]{aone@aoneinst.edu}
\urladdr{http://www.authorone.oneuniv.edu}
\thanks{Thanks for Author One.}
\author{Author Two}
\curraddr[A. Two]{Author Two current address line 1\\
Author Two current address line 2}
\email[A.~Two]{atwo@atwoinst.edu}
\urladdr{http://www.authortwo.twouniv.edu}
\thanks{Thanks for Author Two.}
\author{Author Three}
\address[A. Three]{Author Three address line 1\\
Author Three address line 2}
\urladdr{http://www.authorthree.threeuniv.edu}
\thanks{This paper is in final form and no version of it will be submitted
for publication elsewhere.}
\subjclass[2000]{Primary 05C38, 15A15; Secondary 05A15, 15A18}
\date{July 27, 2001}
\keywords{Keyword one, keyword two, keyword three}
\dedicatory{Dedicated to the memory of S. Bach.}

\begin{abstract}
Replace this text with your own abstract.
\end{abstract}

\maketitle

\section{Symplectic vector space}

In mathematics, a symplectic vector space is a vector space $V$ over a field 
$\mathbb{F}$ (for example the real numbers $\mathbb{R}$) equipped with a
symplectic bilinear form.

A \emph{symplectic bilinear form} is a mapping 
\begin{equation}
\U{3c9} :V\times \ V\rightarrow F
\end{equation}
that is

\begin{itemize}
\item bilinear: linear in each argument separately,

\item alternating: $\U{3c9} (v,v)=0$ holds for all $v\in V$, and

\item nondegenerate: $\U{3c9} (u,v)=0$ for all $v\in V$ implies that $u$ is
zero.
\end{itemize}

If the underlying field has characteristic not $2$, alternation is
equivalent to skew-symmetry. If the characteristic is $2$, the skew-symmetry
is implied by, but does not imply alternation. In this case every symplectic
form is a symmetric form, but not vice versa. Working in a fixed basis, $%
\U{3c9} $ can be represented by a matrix. The conditions above say that this
matrix must be skew-symmetric, nonsingular, and hollow (all diagonal
elements are zero). This is not the same thing as a symplectic matrix, which
represents a symplectic transformation of the space. If $V$ is
finite-dimensional, then its dimension must necessarily be even since every
skew-symmetric, hollow matrix of odd size has determinant zero. Notice the
condition that the matrix be hollow is not redundant if the characteristic
of the field is $2$. A symplectic form behaves quite differently from a
symmetric form, for example, the scalar product on Euclidean vector spaces.

\subsection{Contents}

\begin{enumerate}
\item Standard symplectic space

\item Analogy with complex structures

\item Volume form

\item Symplectic map

\item Symplectic group

\item Subspaces

\item Heisenberg group

\item See also

\item References
\end{enumerate}

\subsection{Standard symplectic space}

The standard symplectic space is $\mathbb{R}^{2n}$ with the symplectic form
given by a nonsingular, skew-symmetric matrix. Typically $\U{3c9} $ is
chosen to be the block matrix%
\begin{equation}
\omega =\left( 
\begin{array}{cc}
\mathbb{O} & -\mathbb{I}_{n} \\ 
\mathbb{I}_{n} & \mathbb{O}%
\end{array}%
\right) 
\end{equation}

where $\mathbb{I}_{n}$ is the $n\times n$ identity matrix. In terms of basis
vectors $\left\{ x_{1},...,x_{n},y_{1},...,y_{n}\right\} $%
\begin{eqnarray}
\omega \left( x_{j},y_{k}\right)  &=&-\omega \left( y_{k},x_{j}\right)
=\delta _{jk} \\
\omega \left( x_{j},x_{k}\right)  &=&\omega \left( y_{j},y_{k}\right) =0 \\
1 &\leq &j,k\leq n
\end{eqnarray}

A modified version of the Gram--Schmidt process shows that any
finite-dimensional symplectic vector space has a basis such that $\U{3c9} $
takes this form, often called a Darboux basis, or symplectic basis.

There is another way to interpret this standard symplectic form. Since the
model space $\mathbb{R}^{2n}$ used above carries much canonical structure
which might easily lead to misinterpretation, we will use "anonymous" vector
spaces instead. Let $V$ be a real vector space of dimension $n$ and $V^{d}$
its dual space. Now consider the direct sum $W=V\oplus V^{d}$ of these
spaces equipped with the following form:%
\begin{equation}
\omega \left( x+u,y+z\right) =z\left( x\right) -y\left( u\right) ;\quad
x,u\in V,\quad y,z\in V^{d}
\end{equation}%
a\TEXTsymbol{\vert}nd consider its dual basis $\left\{ v_{1}^{d},\cdots
,v_{n}^{d}\right\} .$ We can interpret the basis vectors as lying in $W$ if
we write $x_{j}$ $=(v_{j},0)$ and $y_{j}$ $=$ $(0,v_{j}^{d})$. Taken
together, these form a complete basis of $W$, $\left\{
x_{1},...,x_{n},y_{1},...,y_{n}\right\} $.

The form $\U{3c9} $ defined here can be shown to have the same properties as
in the beginning of this section. On the other hand, every symplectic
structure is isomorphic to one of the form $V\oplus V^{d}$. The subspace $V$
is not unique, and a choice of subspace $V$ is called a \emph{polarization}.
The subspaces that give such an isomorphism are called \emph{Lagrangian
subspaces} or simply \emph{Lagrangians}.

Explicitly, given a Lagrangian subspace (as defined below), then a choice of
basis $\left\{ x_{1},...,x_{n}\right\} $ defines a dual basis for a
complement, by%
\begin{equation}
\U{3c9} (x_{j},y_{k})=\U{3b4} _{jk}.
\end{equation}

\subsection{Analogy with complex structures}

Just as every \emph{symplectic structure} is isomorphic to one of the form $%
V\oplus V^{d}$, every \emph{complex structure} on a vector space is
isomorphic to one of the form $V\oplus V$. Using these structures, the \emph{%
tangent bundle} of an $n$-manifold, considered as a $2n$-manifold, has an
almost complex structure, and the \emph{cotangent bundle} of an $n$%
-manifold, considered as a $2n$-manifold, has a symplectic structure: 
\begin{equation}
T_{\ast }(T^{\ast }M)_{p}=T_{p}(M)\oplus (T_{p}(M))^{d}.
\end{equation}

The \emph{complex analog to a Lagrangian subspace is a real subspace }$V$, a
subspace whose complexification is the whole space:%
\begin{equation}
W=V\oplus J_{V}V.
\end{equation}
As can be seen from the standard symplectic form above, every symplectic
form on $\mathbb{R}^{2n}$ is isomorphic to the imaginary part of the
standard complex (Hermitian) inner product on $\mathbb{C}^{n}$ (with the
convention of the first argument being anti-linear).

\subsection{Volume form}

Let $\U{3c9} $ be an alternating bilinear form on an $n$-dimensional real
vector space $V,\U{3c9} \in \wedge ^{2}V$. Then $\U{3c9} $ is non-degenerate
if and only if $n$ is even and 
\begin{equation}
\U{3c9} ^{n/2}=\U{3c9} \wedge ...\wedge \U{3c9} 
\end{equation}
is a volume form. A volume form on a $n$-dimensional vector space $V$ is a
non-zero multiple of the $n$-form $e_{1}^{d}\wedge ...\wedge e_{n}^{d}$
where $\left\{ e_{1},...,e_{n}\right\} $ is a basis of $V$.

For the standard basis defined in the previous section, we have%
\begin{equation}
\U{3c9} ^{n}=(-1)^{n/2}x_{1}^{d}\wedge \ldots \ \wedge x_{n}^{d}\wedge
y_{1}^{d}\wedge \ldots \ \wedge y_{n}^{d},
\end{equation}%
which by reodering becomes 
\begin{equation}
\U{3c9} ^{n}=(-1)^{n/2}x_{1}^{d}\wedge y_{1}^{d}\wedge \ldots \ \wedge
x_{n}^{d}\wedge \wedge y_{n}^{d}\ .
\end{equation}

Authors variously define $\U{3c9} ^{n}$ or $(-1)^{n/2}\U{3c9} ^{n}$ as the
standard volume form. An occasional factor of $n!$ may also appear,
depending on whether the definition of the alternating product contains a
factor of $n!$ or not. The volume form defines an \emph{orientation on the
symplectic vector space} $\left\{ V,\U{3c9} \right\} .$

\subsection{Symplectic map}

Suppose that $\left\{ V,\U{3c9} \right\} $ and $\left\{ W,\U{3c1} \right\} $
are \emph{symplectic vector spaces.} Then a linear map 
\begin{equation}
f:V\rightarrow W
\end{equation}
is called a symplectic map if the pullback preserves the symplectic form,
i.e. $f^{d}\left( \U{3c1} \right) =\U{3c9} $, where the pullback form is
defined by 
\begin{equation}
f^{d}\U{3c1} (u,v)=\U{3c1} (f(u),f(v)).
\end{equation}
\emph{Symplectic maps are volume- and orientation-preserving}.

\subsection{Symplectic group}

If $V=W$, then a symplectic map $f$ is called a linear \emph{symplectic
transformation} of $V$. In particular, in this case one has that 
\begin{equation}
\U{3c9} (f(u),f(v))=\U{3c9} (u,v),
\end{equation}
and so the linear transformation $f$ preserves the symplectic form. The set
of all symplectic transformations forms a group and in particular a Lie
group, called the symplectic group and denoted by $Sp(V)$ or sometimes $Sp(V,%
\U{3c9} )$. In matrix form symplectic transformations are given by
symplectic matrices.

\subsection{Subspaces}

Let $W$ be a linear subspace of $V$. Define the \emph{symplectic complement}
of $W$ to be the subspace%
\begin{equation}
W^{\bot }=\{v\in V\mid \U{3c9} (v,w)=0\,\forall w\in W\}
\end{equation}

The symplectic complement satisfies:

\begin{itemize}
\item 
\begin{equation}
\left( W^{\bot }\right) ^{\bot }=W
\end{equation}

\item 
\begin{equation}
\dim \left\{ W\right\} +\dim \left\{ W^{\bot }\right\} =\dim \left\{
V\right\} 
\end{equation}
\end{itemize}

However, unlike orthogonal complements, $W^{\bot }\cap W$ need not be $\{0\}$%
. We distinguish four cases:

\begin{enumerate}
\item $W$ is \emph{symplectic} if $W^{\bot }\cap W=\{0\}$. This is true if
and only if $\U{3c9} $ restricts to a nondegenerate form on $W$. A
symplectic subspace with the restricted form is a symplectic vector space in
its own right.

\item $W$ is \emph{isotropic} if $W\subseteq W^{\bot }$. This is true if and
only if $\U{3c9} $ restricts to $0$ on $W$. Any one-dimensional subspace is
isotropic.

\item $W$ is \emph{coisotropic} if $W^{\bot }\subseteq W$. $W$ is
coisotropic if and only if $\U{3c9} $ descends to a nondegenerate form on
the quotient space $W/W^{\bot }$. Equivalently $W$ is coisotropic if and
only if $W^{\bot }$ is isotropic. Any codimension-one subspace is
coisotropic.

\item $W$ is \emph{Lagrangian} if $W=W^{\bot }$. A subspace is Lagrangian if
and only if it is both isotropic and coisotropic. In a finite-dimensional
vector space, a Lagrangian subspace is an isotropic one whose dimension is
half that of $V$. Every isotropic subspace can be extended to a Lagrangian
one.
\end{enumerate}

Referring to the canonical vector space $\mathbb{R}^{2}$ above,

\begin{itemize}
\item the subspace spanned by $\{x_{1},y_{1}\}$ is symplectic

\item the subspace spanned by $\{x_{1},x_{2}\}$ is isotropic

\item the subspace spanned by $\left\{ x_{1},...,x_{n},y_{1}\right\} $ is
coisotropic

\item the subspace spanned by $\left\{ x_{1},...,x_{n}\right\} $ is
Lagrangian.
\end{itemize}

\subsection{Heisenberg group}

A \emph{Heisenberg }group can be defined for any symplectic vector space,
and this is the typical way that Heisenberg groups arise.

A vector space can be thought of as a commutative Lie group (under
addition), or equivalently as a commutative Lie algebra, with trivial Lie
bracket. The Heisenberg group is a \emph{central extension} of such a
commutative Lie group/algebra: the symplectic form defines the commutation,
analogously to the \emph{canonical commutation relations (CCR)}, and a \emph{%
Darboux basis} corresponds to canonical coordinates -- in physics terms, to
momentum operators and position operators.

Indeed, \emph{by} the \emph{Stone--von Neumann} theorem, every
representation satisfying the CCR (every representation of the Heisenberg
group) is of this form, or more properly unitarily conjugate to the standard
one.

Further, the group algebra of a vector space is the \emph{symmetric algebra
genrated by that vector space}, and the group algebra of the Heisenberg
group is the \emph{Weyl algebra}: one can think of the \emph{central
extension} as corresponding to \emph{quantization or deformation}.

Formally, the symmetric algebra of $V$ is the group algebra of the dual, 
\begin{equation}
\func{Sym}\left\{ V\right\} =\vee \left( V^{d}\right) \equiv K[V^{d}]
\end{equation}%
and the Weyl algebra is the group algebra of the (dual) Heisenberg group 
\begin{equation}
\func{Weyl}\left\{ V\right\} =K[H(V^{d})]
\end{equation}%
Since passing to group algebras is a \emph{contravariant functor}, the
central extension map 
\begin{equation}
H(V)\rightarrow V
\end{equation}
becomes an inclusion%
\begin{equation}
\func{Sym}\left\{ V\right\} \hookrightarrow W(V).
\end{equation}

\subsection{See also}

\begin{itemize}
\item  A symplectic manifold is a smooth manifold with a smoothly-varying
closed symplectic form on each tangent space

\item  Maslov index

\item A symplectic representation is a group representation where each group
element acts as a symplectic transformation.
\end{itemize}

\end{document}
